\documentclass[12pt,a4paper,twoside]{book}
\usepackage{graphicx}
\usepackage{setspace}	%double spacing for text, single for captions, footnotes, etc.
%\usepackage{hypernat} 	%substitut de cite que permet fer hyperlinks
\usepackage{natbib}		% substituye a 'hypernat' que funciona en Windows.
\usepackage[english]{babel}
\usepackage[utf8]{inputenc}
\usepackage{color}
\usepackage{hhline} 		% extended styles for tables
\usepackage{multirow}
\usepackage{subfigure}
\usepackage{acronym}
\usepackage{hyperref}
\usepackage{amsmath,amsmath,amssymb} 
\usepackage{fancyhdr}
\usepackage{epsfig, amsmath}
\usepackage{algorithm}
\usepackage{algorithmic}

% general settings
\hypersetup{
	linktocpage=true,
	colorlinks=true,
	linkcolor=blue,
	citecolor=blue,
}
\definecolor{Hgray}{gray}{0.6}

\newenvironment{definition}[1][Definition]{\begin{trivlist}
\item[\hskip \labelsep {\bfseries #1}]}{\end{trivlist}}

\setlength{\topmargin}{0cm}
\setlength{\textheight}{23cm}
\setlength{\textwidth}{17cm}
\setlength{\oddsidemargin}{0cm}
\setlength{\evensidemargin}{0cm}
\setlength{\headheight}{1cm}

% indica que las 'sub-sub-sections' sean numeradas y aparezcan en el indice
\setcounter{secnumdepth}{3}
\setcounter{tocdepth}{2}

% settings for code
\renewcommand{\algorithmicrequire}{\textbf{Entrada: }}
\renewcommand{\algorithmicensure}{\textbf{Salida: }}

%%%%%%%%%%%%
% DOCUMENT %
%%%%%%%%%%%%
\begin{document}

% portada
\newpage
\thispagestyle{empty}

\baselineskip 2em

%\vspace*{1cm}

\centerline{\includegraphics[width=0.6\textwidth]{../images/UOC-logo}}
\begin{center}
\textsc{Universitat Oberta de Catalunya (UOC) \\
 Master's Degree in Data Science (\textit{Data Science})\\}

%\centerline {\pic{UOC}{4cm}}

\vspace*{1.5cm}

\textsc{\Large MASTER'S THESIS}

\vspace*{0.5cm}

\textsc{\large Area: Data Mining and Machine Learning}


%\textbf{\Huge VirtualTechLab Model: }

\vspace*{2.0cm}

\textbf{\Large Study of the Transcriptional Function of Cyclin D1 in Leukemia}

\textbf{\large Comparison of gene expression regulated by cyclin D1 in lymphomas with DNA damage.}

\vspace{2.5cm}
\baselineskip 1em

\baselineskip 2em
-----------------------------------------------------------------------------\\
Author:     Antonio Milán Otero\\
Tutor:      Carles Barceló\\
Professor:  Jordi Casas Roma\\
-----------------------------------------------------------------------------\\
\vspace*{1.5cm}
Barcelona, \today

\end{center}

\newpage
\pagestyle{empty}
\hfill

\newpage
% abstract
\pagenumbering{roman} 
\setcounter{page}{1} 
\pagestyle{plain}

%%%%%%%%%%%%%%%%
%%% CREDITOS %%%
%%%%%%%%%%%%%%%%
\chapter*{Créditos/Copyright}

Una página con la especificación de créditos/copyright para el proyecto (ya sea aplicación por un lado y documentación por el otro, o unificadamente), así como la del uso de marcas, productos o servicios de terceros (incluidos códigos fuente). Si una persona diferente al autor colaboró en el proyecto, tiene que quedar explicitada su identidad y qué hizo.

A continuación se ejemplifica el caso más habitual, aunque se puede modificar por cualquier otra alternativa:

\vspace{1cm}

\begin{figure}[ht]
    \centering
	\includegraphics[scale=1]{../images/license.png}
\end{figure}

Esta obra está sujeta a una licencia de Reconocimiento -  NoComercial - SinObraDerivada

\href{https://creativecommons.org/licenses/by-nc-nd/3.0/es/}{3.0 España de CreativeCommons}.

%%%%%%%%%%%%%
%%% FICHA %%%
%%%%%%%%%%%%%
\chapter*{THESIS INDEX CARD}

\begin{table}[ht]
	\centering{}
	\renewcommand{\arraystretch}{2}
	\begin{tabular}{r | p{8cm}}
		\hline
		Title: & Study of the transcriptional function \newline of Cyclin D1 in Leukemia\\
		\hline
        Author: & Antonio Milán Otero\\
		\hline
        Nombre del colaborador/a docente: & Carles Barceló\\
		\hline
        Nombre del PRA (PRA???): & Jordi Casas Roma\\
		\hline
        Date of delivery (mm/aaaa): & 06/2019\\
		\hline
        Degree or program: & MSc in Data Science\\
		\hline
        Thesis area: & Data Mining and Machine Learning\\
		\hline
        Language: & English\\
		\hline
        Keywords & Leukemia, transcriptional function, cyclin D1, machine learning\\
		\hline
	\end{tabular}
\end{table}

%%%%%%%%%%%%%%%%%%%
%%% DEDICATORIA %%%
%%%%%%%%%%%%%%%%%%%
\chapter*{Dedication/Quote}

Breves palabras de dedicatoria y/o una cita.

%%%%%%%%%%%%%%%%%%%
%%% Agradecimientos %%%
%%%%%%%%%%%%%%%%%%%
\chapter*{Acknowledgment}

Si se considera oportuno, mencionar a las personas, empresas o instituciones que hayan contribuido en la realización de este proyecto.

%%%%%%%%%%%%%%%%
%%% RESUMEN  %%%
%%%%%%%%%%%%%%%%
\chapter*{Abstract}
\addcontentsline{toc}{chapter}{Abstract}

\onehalfspacing

Texto con la síntesis del proyecto, esto es, un texto en el cual se explica de manera concisa la definición del proyecto/problema abordado, sus objetivos/métodos de resolución, y los resultados y conclusiones (no puede ser una lista, sino un texto continuo redactado de manera estructurada). Si es necesario poner una referencia en este texto, ésta será anotada a pie de la misma página. En este apartado se puede usar un lenguaje más literario y coloquial que para el resto del documento.

El Abstract se escribirá por duplicado. Una de las versiones tiene que ser \textbf{obligatoriamente en inglés}. La otra versión tiene que estar escrita en catalán o español. En caso de no escribir el resto del documento en inglés, será necesario escribir la segunda versión del Abstract en el idioma utilizado para el resto de la memoria. La palabra Abstract se cambiará por ``\textbf{Resum}'' o ``\textbf{Resumen}'' en la versión catalana y española, respectivamente. 

Extensión recomendada: 250 palabras máximo.

Como escribir un buen Abstract (en inglés):

\href{http://www.ece.cmu.edu/~koopman/essays/abstract.html}{http://www.ece.cmu.edu/~koopman/essays/abstract.html}

\vspace{1.5cm}

\textbf{Palabras clave}: Leukemia, transcriptional function, cyclin D1, machine learning
\newpage

\pagestyle{fancy}
\renewcommand{\chaptermark}[1]{ \markboth{#1}{}}
\renewcommand{\sectionmark}[1]{\markright{ \thesection.\ #1}}
\lhead[\fancyplain{}{\bfseries\thepage}]{\fancyplain{}{\bfseries\rightmark}}
\rhead[\fancyplain{}{\bfseries\leftmark}]{\fancyplain{}{\bfseries\thepage}}
\cfoot{}

% indice
\cleardoublepage
\phantomsection
\addcontentsline{toc}{chapter}{Index}
\tableofcontents
% listado de figuras
\cleardoublepage
\phantomsection
\addcontentsline{toc}{chapter}{List of Figures}
\listoffigures
% listado de tablas
\cleardoublepage
\phantomsection
\addcontentsline{toc}{chapter}{List of Tables}
\listoftables

\thispagestyle{empty}

\pagenumbering{arabic}

\pagestyle{fancy}
\renewcommand{\chaptermark}[1]{ \markboth{#1}{}}
\renewcommand{\sectionmark}[1]{\markright{ \thesection.\ #1}}
\lhead[\fancyplain{}{\bfseries\thepage}]{\fancyplain{}{\bfseries\rightmark}}
\rhead[\fancyplain{}{\bfseries\leftmark}]{\fancyplain{}{\bfseries\thepage}}
\cfoot{}

\onehalfspacing

% capitulos del documento
\chapter{Introduction}
\label{chapter:introduction}


%%% SECTION
\section{General description of the problem}

La leucemia es un conjunto de procesos tumorales de que provoca un aumento descontrolado de leucocitos (glóbulos blancos) en la sangre u órganos linfáticos. La ciclina D1 es un oncogén frecuentemente sobreexpresado en cáncer sobretodo en Leucemia. Es conocido por ser uno de los principales reguladores el ciclo celular aunque se desconoce su función como regulador de la transcripción (proceso por el que se generan las proteínas que controlan todos los procesos celulares). Se sabe que la ciclina D1 se une a las regiones promotoras de muchos genes aunque se desconoce el resultado de su actividad transcripcional que se cree fundamental en el desarrollo de la Leucemia. Recientes estudios de análisis transcripcional global han generado una enorme cantidad de datos susceptibles de ser analizados por data mining y machine learning, siendo su interpretación fundamental para conocer los mecanismos básicos de la célula que pueden conducir a la Leucemia. Evidentemente, la generación de nuevos medicamentos dependerá de conocer bien estos procesos.

\onehalfspacing
% \chapter{Motivation}
% \label{chapter:motivation}

%%% SECTION
\section{Motivation}

\subsection{Why this project?}
The amount of data collected in scientific researches has increased exponentially during the last decades, making the usage of Data Science methodologies a good fit for improving the final analysis and results.
This project is a clear example of how the advances in Data Science can trigger new ways of doing science, expanding the existing tools in order to achieve better results.

\subsection{What can I add?}
During my MSc in Data Science, I have been learning about all the different aspects of a Data Science project, starting from a project management point of view and continuing with all the different phases of acquisition, storage, hypothesis and modeling, visualization and deployment.
For this specific project, though, my focus will be in the area of Data Mining and Machine Learning, and that is what I think I can add to the project, my accumulated experience in the commented area.

\subsection{Personal interest}
My personal interest in this project comes from the fact, or bad luck, of having close family and friends affected for leukemia, therefore, as soon as I saw the proposal of this project I felt emotionally connected to it.

Apart from that first reason, I also consider that one of the best usages of the advances of Data Mining and Machine Learning is to help in the creation of a better society, being one of its strongest foundations the improvement of the quality of the health of each individual. Therefore, I feel responsible for using my new acquired knowledge in areas that can lead to that goal.

\subsection{How can this project improve my CV?}
At the time of writing this, I am working as a research engineer in the control software group in MAX IV Laboratory, a synchrotron located in the south of Sweden that has the purpose of improving the scientific researches in a global encompass.
Until now, my main activities has been related to the build of software for all the different aspects of the control system, from the very low level control, writing drivers for equipment, up to the high level software like graphical user interfaces that enable the scientist to perform their jobs, passing through all the layers in between, like software libraries, servers, etc.
In other words, I have been always close to the control, synchronization and data acquisition, with this project, I can expand my coverage and help also in the next phase of a scientific research, the analysis of the generated data.


\section{Project Objectives}
\subsection{Primary}
\subsection{Secondaries}
\section{Description of the Methodology Used}
\section{Project Research Plan}

% bibliografia
\addcontentsline{toc}{chapter}{Bibliografía}
\bibliographystyle{plain}
\bibliography{referencias}

\end{document}